% Part:sets-functions-relations
% Chapter: size-of-sets
% Section: comparing-sizes
% Hindi translation (hi-Deva-IN) of Open Logic commit 9620cc73f9c8e0ad003c514a5d3748f29611c4c0.

\documentclass[../../../include/open-logic-section]{subfiles}

\begin{document}

\olfileid[hi]{sfr}{siz}{car}

\olsection{भिन्न आकार के समुच्चय और कैंटर का प्रमेय}

\begin{explain}
हमने इस विचार का यथार्थ कथन दिया है कि दो समुच्चयों का आकार समान है। हम इस
विचार का भी यथार्थ कथन दे सकते हैं कि एक समुच्चय दूसरे से छोटा है। ``से छोटा
(या समसंख्यक) है'' की हमारी परिभाषा में समुच्चयों के बीच किसी
!!a{bijection} के स्थान पर पहले समुच्चय से दूसरे में कोई !!a{injection}
चाहिए। यदि ऐसा फलन विद्यमान है, तो पहले समुच्चय का आकार दूसरे समुच्चय के आकार
से कम या उसके बराबर है। सहज रूप में, एक समुच्चय से दूसरे में कोई
!!a{injection} यह सुनिश्चित करता है कि फलन के परिसर में उसके प्रांत के
!!{element} जितने तो !!{element} अवश्य हैं, क्योंकि प्रांत के कोई भी दो
!!{element} परिसर के एक ही अवयव पर नहीं भेजे जाते।
\end{explain}

\begin{defn}
$A$,~$B$ से \emph{बड़ा नहीं है}, जिसे $\cardle{A}{B}$ लिखते हैं, यदि और केवल
यदि कोई !!a{injection} $f \colon A \to B$ हो।
\end{defn}

स्पष्ट है कि यह संबंध स्वतुल्य और संक्रामक है, पर सममित नहीं है (इसे अभ्यास
के रूप में छोड़ा गया है)। हम ऐसी धारणा भी दे सकते हैं जो कहती है कि एक
समुच्चय दूसरे से (वास्तव में) छोटा है।

\begin{defn}
$A$,~$B$ से \emph{छोटा है}, जिसे $\cardless{A}{B}$ लिखते हैं, यदि और केवल
यदि कोई !!a{injection}~$f\colon A \to B$ हो, पर कोई !!{bijection}~$g\colon A
\to B$ न हो; अर्थात $\cardle{A}{B}$ और $\cardneq{A}{B}$।
\end{defn}

स्पष्ट है कि यह संबंध अस्वतुल्य और संक्रामक है। (इसे अभ्यास के रूप में छोड़ा
गया है।) इस संकेतन से हम कह सकते हैं कि कोई समुच्चय $A$ !!{enumerable} है
यदि और केवल यदि $\cardle{A}{\Nat}$, और $A$ !!{nonenumerable} है यदि और केवल
यदि $\cardless{\Nat}{A}$। इससे हम
\oliflabeldef{sfr:siz:nen-alt:thm:nonenum-pownat}{%
\olref[sfr][siz][nen-alt]{thm:nonenum-pownat}
को इस प्रेक्षण के रूप में फिर कह सकते हैं कि
$\cardless{\Nat}{\Pow{\Nat}}$}{%
\olref[sfr][siz][nen]{thm:nonenum-pownat}
को इस प्रेक्षण के रूप में फिर कह सकते हैं कि
$\cardless{\PosInt}{\Pow{\PosInt}}$}। वास्तव में, \citet{Cantor1892} ने
सिद्ध किया कि यह अंतिम बात \emph{पूरी तरह सामान्य} है:

\begin{thm}[कैंटर]\ollabel{thm:cantor}
$\cardless{A}{\Pow{A}}$, किसी भी समुच्चय $A$ के लिए।
\end{thm}

\begin{proof}
प्रतिचित्रण $f(x) = \{x\}$ एक !!a{injection} $f \colon A \to \Pow{A}$
है, क्योंकि यदि $x \neq y$, तो विस्तारिता से $\{x\} \neq \{y\}$ भी, और
इसलिए $f(x) \neq f(y)$। अतः $\cardle{A}{\Pow{A}}$।

\begin{editorial}
यदि \olref[nen]{sec} उपस्थित है, तो हम धीमा प्रमाण देते हैं; अन्यथा
\olref[nen-alt]{sec} से मेल खाने वाला अधिक तीव्र प्रमाण देते हैं।
\end{editorial}

\oliflabeldef{sfr:siz:nen:sec}{% 
  अब हम दिखाएँगे कि कोई !!a{surjective} फलन~$g\colon A \to
  \Pow{A}$ हो ही नहीं सकता, !!a{bijective} फलन तो दूर की बात है; अतः
  $\cardneq{A}{\Pow{A}}$। मान लीजिए कि $g\colon A \to \Pow{A}$। चूँकि
  $g$ सर्वत्र परिभाषित है, प्रत्येक $x \in A$ किसी उपसमुच्चय $g(x)
  \subseteq A$ पर भेजा जाता है। हम दिखा सकते हैं कि $g$ आच्छादक नहीं हो
  सकता। इसके लिए हम एक उपसमुच्चय~$\overline{A} \subseteq A$ परिभाषित करते
  हैं जो परिभाषा के कारण~$g$ के परिसर में नहीं हो सकता। मान लें
  \[
  \overline{A} = \Setabs{x \in A}{x \notin g(x)}.
  \]
  चूँकि $g(x)$ प्रत्येक $x \in A$ के लिए परिभाषित है, इसलिए
  $\overline{A}$ स्पष्टतः~$A$ का सुव्याख्यायित उपसमुच्चय है। पर वह~$g$ के
  परिसर में नहीं हो सकता। कोई भी $x \in A$ लें; हम दिखाएँगे कि $\overline{A} \neq
  g(x)$। यदि $x \in g(x)$, तो वह $x \notin g(x)$ को
  संतुष्ट नहीं करता, इसलिए~$\overline{A}$ की परिभाषा से $x \notin
  \overline{A}$। यदि $x \in \overline{A}$, तो उसे~$\overline{A}$ का
  परिभाषक गुणधर्म संतुष्ट करना होगा; अर्थात $x \in A$ और $x \notin
  g(x)$। चूँकि $x$ कोई भी था, इससे दिखता है कि प्रत्येक $x \in
  A$ के लिए $x \in g(x)$ यदि और केवल यदि $x \notin \overline{A}$, इसलिए
  $g(x) \neq \overline{A}$। दूसरे शब्दों में, $\overline{A}$,~$g$ के
  परिसर में नहीं हो सकता; यह~$g$ के आच्छादक होने की धारणा का खंडन करता
  है।}{अब केवल $\cardneq{A}{\Pow{A}}$ दिखाना शेष है। विरोधाभासार्थ मान
  लीजिए $\cardeq{A}{\Pow{A}}$; अर्थात कोई !!{bijection}
  $g \colon A \to \Pow{A}$ है। अब इस पर विचार करें:
  \[
    D = \Setabs{x \in A}{x \notin g(x)}
  \]
  ध्यान दें कि $D \subseteq A$, इसलिए $D \in \Pow{A}$। चूँकि $g$ एक
  !!a{bijection} है, इसलिए कोई $y \in A$ ऐसा है कि $g(y) = D$। लेकिन
  अब हमारे पास है:
  \[
    y \in g(y) \text{ यदि और केवल यदि } y \in D \text{ यदि और केवल यदि } y \notin g(y).
  \]
  यह विरोधाभास है; अतः $\cardneq{A}{\Pow{A}}$।}{}
\end{proof}

\begin{explain}
\oliflabeldef{sfr:siz:nen:thm:nonenum-pownat}{\olref{thm:cantor} के
  प्रमाण की तुलना \olref[nen]{thm:nonenum-pownat} के प्रमाण से करना
  उपयोगी है। वहाँ हमने दिखाया था कि किसी भी सूची
  $Z_1$, $Z_2$, \dots{}, जो~$\PosInt$ के उपसमुच्चयों की हो, के लिए
  संख्याओं का कोई समुच्चय~$\overline{Z}$ बनाया
  जा सकता है जिसके सूची में न होने की गारंटी है। उसके सूची में न होने की
  गारंटी इसलिए थी कि प्रत्येक $n \in
  \PosInt$ के लिए $n \in Z_n$ यदि और केवल यदि $n \notin \overline{Z}$। इस प्रकार सदैव कोई संख्या होती है जो
  $Z_n$ या $\overline{Z}$ में से एक की !!a{element} है, पर दूसरे की नहीं।
  यहाँ हम इसी विचार का अनुसरण करते हैं, अंतर केवल यह है कि अब सूचकांक~$n$,
  ~$A$ के !!{element},~$\PosInt$ के बजाय हैं। समुच्चय $\overline{A}$ इस
  प्रकार परिभाषित है कि वह~$g(x)$ से प्रत्येक $x \in A$ के लिए भिन्न है,
  क्योंकि $x \in g(x)$ यदि और केवल यदि $x \notin \overline{A}$। फिर से,
  ~$A$ का सदैव कोई !!a{element} होता है जो $g(x)$ और $\overline{A}$ में से
  एक का !!a{element} है, पर दूसरे का नहीं। और जैसे इस कारण $\overline{Z}$
  सूची $Z_1$, $Z_2$, \dots{} में नहीं हो सकता, वैसे ही $\overline{A}$,
  ~$g$ के परिसर में नहीं हो सकता।}{}

\oliflabeldef{sfr:siz:nen-alt:thm:nonenum-pownat}{\olref{thm:cantor} के
  प्रमाण की तुलना \olref[nen-alt]{thm:nonenum-pownat} के प्रमाण से करना
  उपयोगी है। वहाँ हमने दिखाया था कि किसी भी सूची
  $N_0$, $N_1$, $N_2$, \dots{}, जो~$\Nat$ के उपसमुच्चयों की हो, के लिए हम
  संख्याओं का कोई समुच्चय~$D$ बना
  सकते हैं जिसके सूची में न होने की गारंटी है। उसके सूची में न होने की
  गारंटी इसलिए थी कि $n \in N_n$ यदि और केवल यदि $n \notin
  D$, प्रत्येक $n \in \Nat$ के लिए। यहाँ हम इसी विचार का अनुसरण करते हैं, अंतर केवल यह है कि
  अब सूचकांक~$n$,~$A$ के !!{element},~$\Nat$ के बजाय हैं। समुच्चय $D$ इस
  प्रकार परिभाषित है कि वह~$g(x)$ से प्रत्येक $x
  \in A$ के लिए भिन्न है,
  क्योंकि $x \in g(x)$ यदि और केवल यदि $x \notin D$।}{}

इस प्रमाण की तुलना रसेल के विरोधाभास के प्रमाण,
\olref[sfr][set][rus]{thm:russells-paradox}, से करना भी उपयोगी है। वास्तव
में कैंटर का प्रमेय ही रसेल के अपने विरोधाभास की प्रेरणा था।
\end{explain}

\begin{prob}
  दिखाइए कि कोई !!a{injection} $g\colon \Pow{A} \to
  A$ किसी भी समुच्चय~$A$ के लिए हो ही नहीं सकता। संकेत: मान लीजिए
  $g\colon \Pow{A} \to A$ !!{injective} है।
  $D = \Setabs{g(B)}{B \subseteq A \text{ और
  } g(B) \notin B}$ पर विचार कीजिए। मान लें $x = g(D)$। इस तथ्य का उपयोग
  करके विरोधाभास निकालिए कि $g$ !!{injective} है।
\end{prob}

\end{document}
